\documentclass[a4paper,10pt,english]{article}

\usepackage[margin=2.5cm]{geometry}
\usepackage[utf8]{inputenc}
\usepackage[T1]{fontenc}
\usepackage[english]{babel}
\usepackage{amsmath}
\usepackage{amssymb}
\usepackage{amsthm}
\usepackage{booktabs}
\usepackage{longtable}
\usepackage{array}
\usepackage[colorlinks,urlcolor=blue,citecolor=blue,linkcolor=blue]{hyperref}

\theoremstyle{definition}
\newtheorem{proposition}{Proposition}

\newcommand{\Prob}{\mathbb P}

\setlength{\emergencystretch}{3em}

\title{Correction Roadmap for Coordinate Search}
\author{Th\'eo Voldoire and Nic Fishman}
\date{}

\begin{document}

\maketitle

\begin{abstract}
This note is a numbered crosswalk from the coordinate-search results in
\emph{The Geometry of p-hacking} to their replacements in
\emph{Unified Moment Approximation for Specification Search}.  It records
what remains valid, what fails, what the unified results recover, and
which additional conditions are necessary.
\end{abstract}

Write \emph{Main} for the first document and \emph{Unified} for the
second.  Result and equation numbers refer to their current numbered
versions.  The scope ends with Main Proposition~1; Main Theorems~6--7
and Appendix~C concern basis search and are unchanged.  All notation is
that of Unified, in particular $m=s\log(ep/s)$ and
$\delta=(m/n)^{1/2}$.

The common replacement spine is Unified Lemma~1.1, which gives the exact
Schur-complement representation; Theorem~2.1, which gives the critical
empirical expansion; Theorem~2.2, which gives the covariance-matched
Gaussian field when $m=o(n)$; and Theorem~2.6, which extends the same
argument to residual-based statistics.  The table below identifies the
endpoint result used for each Main claim.

\section{Numbered crosswalk}

\begingroup
\small
\setlength{\tabcolsep}{3.2pt}
\renewcommand{\arraystretch}{1.06}
\begin{longtable}{@{}>{\raggedright\arraybackslash}p{0.16\textwidth}
>{\raggedright\arraybackslash}p{0.28\textwidth}
>{\raggedright\arraybackslash}p{0.29\textwidth}
>{\raggedright\arraybackslash}p{0.21\textwidth}@{}}
\toprule
Main result & Disposition and exact defect & Unified replacement &
Net change and refactor action \\
\midrule
\endfirsthead
\toprule
Main result & Disposition and exact defect & Unified replacement &
Net change and refactor action \\
\midrule
\endhead

Theorem~1, upper half of (12) &
\emph{Correct.}  The experimental coefficient upper rate and its
probability are valid. &
Theorem~3.1, equations (54)--(55), specialized through (61). &
Keep the conclusion and replace its proof.  The rate is unchanged,
Main Assumption~2 is unnecessary, and the bound is uniform over
$|S|\le s$. \\
\addlinespace

Theorem~1, lower half of (12); Appendix Lemma~7, (35) &
\emph{False without multiplicity.}  The norm $\Psi_s$ can be carried by
one support and does not supply a $\sqrt m$ search gain.  Lemma~7's
inverse-Gram greedy transfer is also false; see
Proposition~\ref{prop:obstructions}(i),(iv). &
Theorem~3.2, (57)--(58), gives the general lower transfer.
Theorem~3.4, (65)--(68), gives the exact scalar replacement;
Lemma~3.5 and Corollary~3.6 give the top-$s$ score-pool construction. &
Replace the lower half.  The old scalar rate is recovered under an
exponentially large separated packing and multiplier minoration.  This
is genuine additional structure; Gaussian treatment supplies the
minoration. \\
\addlinespace

Theorem~2, (14); Appendix Lemma~10, (37) &
\emph{False as written.}  The bound omits the treatment-projection
fluctuation already visible in Lemma~10; see
Proposition~\ref{prop:obstructions}(ii). &
Theorem~3.1, (54)--(55), gives the corrected scale
$\frac{\sigma_r}{\sigma_0}\delta
(\sqrt{\phi_s}+\sqrt{\theta_s^r}+\delta)$. &
Replace the theorem.  The bound can be larger by the genuine
$\frac{\sigma_r}{\sigma_0}\sqrt{\phi_s}\,\delta$ term, but it drops
restricted conditioning, permits singular control blocks, and needs
only sparse overlap.  The old rate returns when this term is absorbed;
Unified (74) is a sufficient condition. \\
\addlinespace

Theorem~3, (17) &
\emph{False.}  The log-cardinality of population-near-optimal supports
is not the canonical entropy of their fluctuations; many such supports
may index the same random variable; see
Proposition~\ref{prop:obstructions}(iii). &
Theorem~3.2, (57)--(58), transfers matched-field quantiles, and
Theorem~3.3, (59)--(60), gives the canonical-metric packing lower bound. &
Replace the theorem.  No unconditional scalar rate follows from
near-optimal cardinality.  Metric separation is additional; conditional
on it, the replacement is covariance-adaptive.  The Gaussian transfer
uses $m=o(n)$. \\
\addlinespace

Theorem~4, upper half of (19); Appendix Lemma~14, (72) &
\emph{False in general.}  The rate omits the treatment channel, the
statement lacks the treatment floor needed by a valid denominator
argument, and the proof invokes hypotheses absent from the theorem.
Its centered studentizer transfer is not proved; see
Proposition~\ref{prop:obstructions}(ii). &
Theorem~5.1: corrected upper bound (100), aligned upper bound (101), and
least-squares transfer (103)--(104).  Corollary~5.2, (109), gives the
matched stochastic approximation. &
Replace the theorem's upper half and its studentization argument.  The
general rate gains $\sqrt m\sqrt{\phi_s}$ and an explicit treatment
floor.  Sparse treatment alignment recovers the old rate.  Main
Assumption~2 is unnecessary for this upper bound. \\
\addlinespace

Theorem~4, lower half of (19) &
\emph{False without direction multiplicity.}  Explained energy
$\theta_s$ does not separate the normalized residual directions that
index the multiplier process; see
Proposition~\ref{prop:obstructions}(i). &
Corollary~5.3, (110), gives matched-field transfer.  Corollary~5.4,
(112)--(115), gives direction-packing bounds; Lemmas~5.5--5.6,
(117)--(123), and Corollary~5.7 give the score-to-direction bridge. &
Replace the lower half.  The old rate becomes conditional on direction
packing, a radius bound, and multiplier minoration.  The direct
critical-regime result is randomized; the general transfer uses
$m=o(n)$. \\
\addlinespace

Theorem~5, (21); Appendix Lemmas~19--20 &
\emph{Not established under Main Assumption~6.}  Entrywise fourth-
cumulant decay does not control sparse score-covariance blocks in
operator norm when $s$ grows, and marginal tail matching does not give
the required joint top-$s$ process. &
Theorem~2.2, (20)--(21), is the general stochastic approximation.
Theorems~4.1--4.3, (75)--(77), (83), and (85), give covariance-aware and
observable top-$s$ reductions.  Theorem~4.4, (91), and Corollary~4.5,
(92), give independent products. &
Replace the theorem by this hierarchy.  The observable empirical sort
is stronger computationally.  Independent products additionally need
score isotropy (89), the mean cap (90), $m\to\infty$, and $m=o(n)$.
The route needs neither $m^4=o(n)$ nor $\log p=o(n^{1/3})$. \\
\addlinespace

Appendix Proposition~1, (78) &
\emph{Not recovered, but not shown false.}  The proof has neither an
approximation with $o_{\Prob}(n^{-1})$ error nor the joint point-process
control needed for its Poisson/Gumbel limit. &
None.  Corollary~4.5, (92), is only a relative-scale approximation and
cannot imply this endpoint. &
Quarantine Proposition~1 as a separate result to be reproved, or remove
it from the refactor. \\

\bottomrule
\end{longtable}
\endgroup

\section{Refactor map}

\begin{center}
\small
\begin{tabular}{@{}>{\raggedright\arraybackslash}p{0.35\textwidth}
>{\raggedright\arraybackslash}p{0.57\textwidth}@{}}
\toprule
Main material & Replacement \\
\midrule
Appendix~A.1--A.3 common machinery &
Unified Lemma~1.1 and Theorems~2.1, 2.2, and 2.6. \\
\addlinespace
Theorems~1--3 and their proofs &
Unified Section~3, especially Theorems~3.1--3.4 and Corollary~3.6. \\
\addlinespace
Theorem~4 and Appendix~A.3 &
Unified Section~5, especially Theorem~5.1 and
Corollaries~5.2--5.4 and 5.7, and Lemmas~5.5--5.6. \\
\addlinespace
Theorem~5 and Appendix~B &
Unified Section~4, Theorems~4.1--4.4 and Corollary~4.5. \\
\addlinespace
Theorems~6--7 and Appendix~C &
No change; these basis-search results lie outside this refactor. \\
\bottomrule
\end{tabular}
\end{center}

The old proof chain should not be ported piecemeal.  Main Lemma~6 leaves
the exceedance concentration open; Lemma~7 is false; Lemma~10 exposes
the term omitted from Theorem~2; Lemma~14 uses assumptions missing from
Theorem~4; and Lemmas~19--20 do not prove a joint extreme-value
comparison.  Appendix~B's fixed-factor extension also leaves Lemma~22
explicitly unproved.  The unified spine replaces their roles directly.

Three numbered introduction displays should be regenerated only after
the theorem refactor.  Main (4) should use Unified Theorem~3.1 and state
matching only through Theorem~3.4 or Corollary~3.6; Main (5) should use
the corrected scale in Unified (54)--(55); and Main (6) should present
the hierarchy in Unified Theorems~4.1--4.4 and Corollary~4.5.  Its
remainder should agree with Unified (92), not claim the stronger
$o_{\Prob}\{s\sqrt{\log(ep/s)}/n\}$ scale.  The condition-number
sentence in Main Section~2.3 should also be deleted or compare matrices
after the same diagonal normalization; interlacing does not compare a
correlation-normalized sparse block with unnormalized $C$.

\appendix
\section{Obstruction certificates}

The following elementary Gaussian examples are retained here because
neither Main nor Unified contains them.  They justify deleting the false
claims rather than merely replacing their proofs.

\begin{proposition}[Four obstructions]\label{prop:obstructions}
Assume that all displayed nonzero variances and coefficients are fixed
and bounded away from zero.
\begin{enumerate}
\item[(i)] \emph{Signal without multiplicity.}
Let $s=1$, $C=I_p$, and let
$x_0,X_1,\ldots,X_p,\varepsilon$ be mutually independent centered
Gaussian variables.  Set
\[
 y=bX_1+\varepsilon,
 \qquad L=\log(ep)\longrightarrow\infty,
 \qquad L=o(\sqrt n).
\]
Then $\Psi_1=|b|$, but
\[
 \max_{j\le p}(\hat\beta_{0,\{j\}}-\hat\beta_{0,\emptyset})
 =O_{\Prob}\left(\frac{|b|}{\sqrt n}
 +\frac{\varsigma L}{n}\right)
 =o_{\Prob}\left(|b|\sqrt{\frac Ln}\right),
 \qquad
 \max_{j\le p}\hat t_{0,\{j\}}-\hat t_{0,\emptyset}=O_{\Prob}(1).
\]
Main Theorems~1 and 4 instead assert lower orders
$|b|\sqrt{L/n}$ and $\sqrt L$.

\item[(ii)] \emph{Diffuse treatment projection.}
Let $p=2s$, $s\to\infty$, $s=o(\sqrt n)$, and let
$X_1,\ldots,X_p,U$ be independent standard Gaussians.  For fixed
$0<\rho<1$, set
\[
 x_0=\frac{\rho\sigma_0}{\sqrt p}\sum_{j=1}^pX_j
 +\sigma_0\sqrt{1-\rho^2}\,U,
 \qquad y=\varepsilon,
\]
with independent Gaussian noise.  Here $\Psi_s=0$ and
$\phi_s=\rho^2/2$, while for any fixed population maximizer $S^\star$,
\[
 \max_{|S|=s}\hat\beta_{0,S}-\hat\beta_{0,S^\star}
 =\Theta_{\Prob}\left(
 \frac{\sigma_\varepsilon}{\sigma_0}\sqrt{\frac sn}\right),
 \qquad
 \max_{|S|=s}\hat t_{0,S}-\hat t_{0,\emptyset}
 =\Theta_{\Prob}(\sqrt s).
\]
Main Theorems~2 and 4 give only
$(\sigma_\varepsilon/\sigma_0)s/n$ and $1+s/\sqrt n$.

\item[(iii)] \emph{Cardinality without fluctuation entropy.}
Let $p=2s$, $s\to\infty$, $s=o(n)$, $C=I_p$, and
\[
 x_0=\rho\sigma_0X_1+\sigma_0\sqrt{1-\rho^2}\,U,
 \qquad y=\beta_0x_0+bX_1,
 \qquad 0<\rho<1,\quad \rho b<0.
\]
Every support containing coordinate $1$ is population-optimal and
returns $\hat\beta_{0,S}=\beta_0$ exactly.  With probability tending to
one, no support omitting coordinate $1$ reaches $\beta_0$.  Thus the
searched excess is zero although
$\ell_0=\log\binom{p-1}{s-1}\asymp s$, contradicting Main Theorem~3.

\item[(iv)] \emph{Inverse-Gram greedy transfer.}
For fixed $0<\rho<1$, let
\[
 R=\begin{pmatrix}1&\rho\\ \rho&1\end{pmatrix},
 \qquad a=(M,1)^\top,
 \qquad b=(1,M)^\top.
\]
Both coordinate products are positive and
$\lambda_{\min}(R)=1-\rho$, but
\[
 a^\top R^{-1}b
 =\frac{2M-\rho(M^2+1)}{1-\rho^2}<0
\]
for large $M$.  A block eigenvalue floor cannot justify Main
Appendix Lemma~7.
\end{enumerate}
\end{proposition}

\begin{proof}
For (i), Unified Theorem~4.1 reduces the coefficient field to score
products: the unique signal coordinate contributes
$O_{\Prob}(|b|/\sqrt n)$, while the largest null products contribute
$O_{\Prob}(L/n)$.  The exact normalized-residual identity in Unified
(116) similarly gives $O_{\Prob}(1+L/\sqrt n)=O_{\Prob}(1)$
for the studentized search.

For (ii), every size-$s$ support has residual treatment variance
$\eta=\sigma_0^2(1-\rho^2/2)$.  The support-varying anchored linear
coefficient field is
\[
 -\eta^{-1}d_S^\top\mathbf X_S^\top
 \boldsymbol\varepsilon/n,
 \qquad d_j=\rho\sigma_0/\sqrt p.
\]
Conditional on $\boldsymbol\varepsilon$, the $p$ scores are independent
Gaussians, and the sum of their largest signed half is
$\Theta_{\Prob}(s)$.  This gives the coefficient order above; the same
sum in the Schur-correlation field, multiplied by $\sqrt n$, gives the
studentized order.

For (iii), a regression containing $X_1$ represents
$y=\beta_0x_0+bX_1$ exactly.  A support omitting $X_1$ has a fixed
negative population gap, which Unified Theorem~3.1 preserves uniformly
with probability tending to one.  Counting the optimal supports gives
the stated $\ell_0$.  Item (iv) is direct matrix inversion.
\end{proof}

\end{document}
